\subsection{A zero-Euler filling which does not admit the smooth-model comparison}

\begin{statement}{C.2.UN.3}{Proved}{The same punctured torus family has non-bimeromorphic zero-Euler fillings}
Let $A=E\times E$, $E=\C/(\Z+\mathrm i\Z)$.
There are two proper smooth-total-space fillings of
$A\times\Delta^*$ over a disc whose reduced central
supports both have Euler number zero, but for which
no bimeromorphic map over the disc exists. One is
$A\times\Delta$. The other is the local filling
(D.23), whose central support consists of a primary
Kodaira surface and a primary Hopf surface meeting
along one elliptic curve. The nonexistence persists
after every finite base change and resolution of the
normalized base change, with the given punctured
family retained.

These fillings occur in compact smooth threefolds:
the principal $A$-bundle $Y_0$ and the threefold $Y$
constructed in D.2.7 are isomorphic away from the
three specified fibres. That isomorphism does not
extend bimeromorphically, even after a finite base
change of their base curve. In particular smooth
torus fillings, together with modifications by
blowups and blowdowns over the central fibre, do
not generate all zero-Euler fillings of this
punctured family.

More strongly, the same punctured product has
infinitely many pairwise non-bimeromorphic
zero-Euler fillings. They are obtained by replacing
$\OO_V(R)$ in (C.UC.2) by $\OO_V(nR)$,
$n\ge1$. This distinction also persists after
arbitrary finite base change and resolution.
\end{statement}
\begin{proof}
We first record the precise birational comparison
used in the proof. Let $\phi:M\dashrightarrow N$
be a bimeromorphic map between smooth complex
threefolds, and let $P\subset M$ be a compact
irreducible divisor whose smooth resolution is
not covered by rational curves. Then $\phi$
cannot contract $P$ to a subset of dimension
at most one. Resolve the inverse meromorphic
map by smooth-centre blowups:
\[
 \widehat N\xrightarrow{\ a\ }N,
 \qquad b:\widehat N\longrightarrow M,
 \qquad b=\phi^{-1}a.
\tag{C.UC.1}
\]
The elimination-of-indeterminacy construction
has this form in the analytic category as
well as the algebraic one; see
\cite[Section~1.2.4 and Lemma~1.3.1]{AKMW2002}.
For maps of proper disc models, the same
construction is made on a smaller disc; its
centres above the compact central fibre are
compact. For the compact threefolds in the
statement the cited construction applies
directly.

Since $b$ is a proper modification of the
smooth target $M$, the strict transform
$\widehat P$ maps birationally onto $P$.
Here is the analytic justification, without
using a global function field. The locus in
$M$ over which $b$ has positive-dimensional
fibres is analytic by properness. It cannot
contain a divisor: its preimage would then
have dimension at least three while being
a proper analytic subset of the irreducible
threefold $\widehat N$. Thus $b$ is finite
over a dense open subset of $P$. A finite
modification of the normal manifold $M$
is an isomorphism, so its degree there is
one. If $\phi$ contracted
$P$, then $a$ would contract $\widehat P$.
It would be an exceptional divisor of the
composition of smooth-centre blowups $a$.
Every such divisor is birational to the
exceptional divisor created at one of those
steps. In dimension three that divisor is
$\PP^2$ for a point centre, or a
$\PP^1$-bundle over a smooth curve for a
curve centre. A divisor centre gives an
isomorphism and creates no exceptional
divisor. Hence $\widehat P$, and therefore
$P$, would be covered by rational curves,
a contradiction. Applying the same argument
to the inverse shows that the surviving
divisors correspond birationally, rather
than only by a generically finite map.

We now give the actual local models. Put
$q=e^{-2\pi}$ and
\[
 V=\operatorname{Bl}_{(0,0)}(E\times\Delta),
 \qquad R=\operatorname{Exc}(V\to E\times\Delta),
 \qquad U=\OO_V(R)^\times/q^{\Z}.
\tag{C.UC.2}
\]
The quotient map $U\to V$ is a principal
elliptic bundle with fibre $\C^*/q^{\Z}=E$.
It is smooth and proper, since each bundle
chart is the product with this compact
elliptic curve. Composition with $V\to\Delta$
is therefore a proper map from a smooth
threefold. The canonical section of
$\OO_V(R)$ trivializes it on the punctured
disc, giving the specified isomorphism
\[
                         U|_{\Delta^*}=A\times\Delta^*.
\tag{C.UC.3}
\]
This is exactly the construction in D.2.7,
including its holomorphic gluing to the
compact threefold.

The strict transform $E_0$ of $E\times\{0\}$
and the exceptional $R$ meet transversely at
one point. The restrictions of $\OO_V(R)$
to them have degree $+1$ and $-1$,
respectively. Their inverse images in $U$
are
\[
 P=\OO_E(0)^\times/q^{\Z},\qquad
 Q=\OO_{\PP^1}(-1)^\times/q^{\Z}
       =(\C^2\setminus\{0\})/\langle q\rangle.
\tag{C.UC.4}
\]
They meet in a single elliptic fibre $J$,
and the local equation of the map is
$xy=t$ along $J$. The first surface is a
primary Kodaira surface with $b_1(P)=3$:
in its principal elliptic bundle the primitive
Chern class kills one of the two elliptic
fibre cycles. The second is the indicated
primary Hopf surface, with $b_1(Q)=1$.
These identifications and the integral
transgression are also calculated in
(D.23)--(D.25).

Neither surface contains a rational curve.
A holomorphic map $\PP^1\to P$ projects
constantly to its elliptic base and then
constantly to its elliptic fibre. A map
$\PP^1\to Q$ lifts to $\C^2\setminus\{0\}$,
since $\PP^1$ is simply connected; its two
holomorphic coordinate functions are constant.
Thus their resolutions are not covered by
rational curves. A bimeromorphic comparison
from $U$ to $A\times\Delta$ over $\Delta$
would, by (C.UC.1), have to carry each of
$P,Q$ birationally onto a central divisor.
The product has only the one central
divisor $A$, and neither $P$ nor $Q$ is
birational to it. Indeed the first Betti
number of a smooth compact surface is
unchanged by point blowups: a ball is
replaced by the simply connected exceptional
neighbourhood, and van Kampen gives the
same fundamental group. Resolution of a
surface bimeromorphism therefore preserves
$b_1$, whereas $b_1(A)=4$ differs from
both values in (C.UC.4). This proves that
there is no such bimeromorphic comparison.

All these supports have Euler number zero.
For $P$ and $Q$, their principal elliptic
bundle projections are locally trivial
proper bundles, so their Euler numbers are
the product of the Euler numbers of base
and fibre, hence zero. The intersection
$J$ also has Euler number zero. Additivity
gives
\[
                    e(P\cup Q)=e(P)+e(Q)-e(J)=0=e(A).
\tag{C.UC.5}
\]

Here is the comparison after a finite base
change, with its ramification retained.
Locally such a change has equation
$t=u(s)s^m$, $u(0)\ne0$. Taking a local
$m$th root of the unit replaces its
parameter by one with $t=s^m$.
At a general point of either $P$ or $Q$
the original map is a submersion with
coordinate $x=t$. Its base change is the
smooth chart
\[
                          x=s^m,
\tag{C.UC.6}
\]
with the two surface coordinates unchanged.
The component over $s=0$ thus maps with
degree one to the original component.
At their intersection the normalized base
change has equation $xy=s^m$, with one
additional smooth elliptic coordinate.
This is normal, with transverse rational
double point $A_{m-1}$ when $m>1$.
A resolution adds ruled exceptional
surfaces above $J$; the strict transforms
of $P$ and $Q$ remain birational to the
surfaces in (C.UC.4). No choice of further
resolution changes those birational
classes. The target base change remains
the smooth product $A\times\Delta_s$.
The argument using (C.UC.1) and their
first Betti numbers again excludes every
bimeromorphic comparison over $\Delta_s$.

Finally use the compact gluing of D.2.7.
The original principal $A$-bundle $Y_0$
has a smooth abelian central fibre at
each of the three chosen base points.
The replacement $Y$ has the local model
(C.UC.2) at each point and is identified
with $Y_0$ elsewhere by (C.UC.3).
An extension of this isomorphism would
respect the base, since the two base
maps agree on its dense open graph.
Its restriction near any of the three
fibres would be the already excluded
comparison. The same argument applies
above every preimage of these points
under a finite cover of the compact base,
using (C.UC.6). This proves the compact
and base-change assertions.

For the infinite family put
\[
 U_n=\OO_V(nR)^\times/q^{\Z},\qquad
 P_n=\OO_E(n[0])^\times/q^{\Z},\qquad
 Q_n=\OO_{\PP^1}(-n)^\times/q^{\Z}.
\tag{C.UC.7}
\]
The canonical section again trivializes the
bundle off the central fibre. Thus the punctured
isomorphism is the same exact product (C.UC.3).
The central support is $P_n\cup_J Q_n$;
all three terms have Euler number zero by their
principal elliptic bundle projections. Each
$P_n$ contains no rational curve by its
projection to $E$ and then to its elliptic
fibre. For $Q_n$, the usual quotient
$\OO_{\PP^1}(-n)^\times
=(\C^2\setminus\{0\})/\boldsymbol\mu_n$
is explicit: a nonzero vector $v$ represents
$v^{\otimes n}$ in the $n$th tensor power of
the tautological line. Multiplication by $q$
on that tensor power lifts to multiplication
by the positive real number $q^{1/n}$ on
$\C^2\setminus\{0\}$. Hence
\[
 Q_n=(\C^2\setminus\{0\})/
       \langle\boldsymbol\mu_n,q^{1/n}\rangle,
 \qquad
 \pi_1(Q_n)=\Z\oplus\Z/n\Z.
\tag{C.UC.8}
\]
The two commuting generators are independent,
because a nonzero power of $q^{1/n}$ has
absolute value different from one. The action
is free and properly discontinuous, and its
covering space retracts to the simply connected
$S^3$. This proves the fundamental-group
calculation and also proves that every map
$\PP^1\to Q_n$ is constant by lifting it to
the cover.

For completeness the corresponding calculation
for $P_n$ retains the extension class. A Hermitian
metric on $\OO_E(n[0])$ identifies its nonzero
total space smoothly with its unit-circle bundle
times $\R_{>0}$. As $q$ is positive real, its
quotient is the unit-circle bundle times
$\R_{>0}/q^{\Z}\simeq S^1$. The circle bundle
has Chern number $n$. Lifting the two generators
of $\pi_1(E)$, the attaching map of the base
two-cell has winding number $n$ in the circle
fibre. With a consistent orientation this gives
the exact presentation
\[
 \pi_1(P_n)=
 \langle a,b,c,d\mid
 [a,b]=c^n,\ [a,c]=[b,c]=1,
 [a,d]=[b,d]=[c,d]=1\rangle.
\tag{C.UC.9}
\]
There are no further relations: the homotopy
sequence of the circle bundle has injective
fibre group because $\pi_2(E)=0$, and the last
factor is the displayed direct $S^1$ factor.
Abelianizing (C.UC.8)--(C.UC.9) yields
\[
 H_1(P_n,\Z)=\Z^3\oplus\Z/n\Z,
 \qquad H_1(Q_n,\Z)=\Z\oplus\Z/n\Z.
\tag{C.UC.10}
\]
Integral first homology, including its torsion,
is invariant under a bimeromorphism of smooth
compact surfaces, by the same point-blowup
comparison already used above. By (C.UC.1),
a bimeromorphism $U_n\dashrightarrow U_{n'}$
over the disc must carry the two non-uniruled
central divisors birationally to the two such
divisors on the target. The ranks in (C.UC.10)
force $P_n$ to correspond to $P_{n'}$ and
$Q_n$ to $Q_{n'}$. The torsion orders then
force $n=n'$. After any finite base change,
(C.UC.6) still leaves these components
birationally unchanged, and the resolution of
$xy=s^m$ adds only ruled surfaces over $J$.
Further resolutions add only uniruled
exceptional divisors. They cannot replace a
non-uniruled component under (C.UC.1), so
the same torsion argument proves persistence.
Finally all $U_n$ have compact realizations:
replace $V$ by
$\operatorname{Bl}_{(0,0)}(E\times\PP^1)$
and use the same line bundle $\OO(nR)$.
Its elliptic quotient is a compact smooth
threefold, and the projection to $\PP^1$
restricts to (C.UC.7) near zero.
\end{proof}

The comparison obstruction specifies the limit
of a particular proposed classification route.
In \cite[Sections~5.1--5.3]{KKN2008}, a
nondegenerate logarithmic torus pairing is
given first, and a suitable complete fan
produces a proper model of that logarithmic
complex torus. The theorem there supplies
those models; an extension of an arbitrary
punctured-family isomorphism to a
bimeromorphic map of its chosen fillings
is an additional requirement. Statement
C.2.UN.3 refutes that additional requirement
for the displayed fillings, even though
the punctured family is a product.
